\section{Заключение}
\label{sec:Chapter7} \index{Chapter7}

%Здесь надо перечислить все результаты, полученные в ходе работы. Из текста
%должно быть понятно, в какой мере решена поставленная задача.

В статье рассматривались криптосистемы $\LAC\,256$ и $\NewHope\,1024$. Процесс шифрования в них был рассмотрен, как передача данных по зашумлённому $\mathsf{RLWE}$-каналу.

Были исследованы свойства этого канала, в предположении независимости компонент шума оценена вероятность ошибочной декодировки сообщения, проанализировано кол-во ошибок, которое необходимо уметь исправлять для получения заданного $\mathsf{DFR}$ в метриках Хэмминга и Ли.

Используя границы Варшамова-Гилберта и Хэмминга, были получены оценки на скорость оптимальной передачи данных по каналу при фиксированной метрике и размере алфавита $Q$.

Получено, что для $\LAC$ оптимальными $Q$ являются $2$ и $3$. При этом метрики Хэмминга и Ли с такими размерами алфавита эквивалентны.
\begin{table}[h]
	\centering
	\begin{minipage}{0.48\textwidth}
		\caption{LAC, Хэмминг, Табл~\ref{tab:LAC,Ham}}
		\centering
		\begin{tabular}{|c|c|c|}
			\hline Q & $R_{GV}$ & $R_H$ \\ \hline
			2 & 0.7569 & 0.8591 \\ \hline
			3 & 0.6298 & 0.9939 \\ \hline
		\end{tabular}
	\end{minipage}
	\hfill
	\begin{minipage}{0.48\textwidth}
		\caption{LAC, Ли, Табл~\ref{tab:LAC,Li}}
		\centering
		\begin{tabular}{|c|c|c|}
			\hline Q & $R_{GV}$ & $R_H$ \\ \hline
			2 & 0.7569 & 0.8591 \\ \hline
			3 & 0.6298 & 0.9939 \\ \hline
		\end{tabular}
	\end{minipage}
\end{table}

Для $\NewHope$ ситуация интереснее, рассмотрим несколько $Q$:
\begin{table}[h]
	\centering
	\begin{minipage}{0.48\textwidth}
		\centering
		\caption{NewHope, Хэмминг, Табл~\ref{tab:Hope,Ham}}
		\begin{tabular}{|c|c|c|}
			\hline Q & $R_{GV}$ & $R_H$ \\ \hline
			8 & 2.9342 & 2.9650 \\ \hline
			14 & 3.1555 & 3.4450 \\ \hline
			34 & 0.4060 & 2.1680 \\ \hline
		\end{tabular}
	\end{minipage}
	\hfill
	\begin{minipage}{0.48\textwidth}
		\centering
		\caption{NewHope, Ли, Табл~\ref{tab:Hope,Li}}
		\begin{tabular}{|c|c|c|}
			\hline Q & $R_{GV}$ & $R_H$ \\ \hline
			8 & 2.9448 & 2.9703 \\ \hline
			14 & 3.3518 & 3.5442 \\ \hline
			34 & 2.8495 & 3.5972 \\ \hline
		\end{tabular}
	\end{minipage}
\end{table}

При маленьких $Q$ ошибок почти не происходит, исправлять их легко, и скорость получается близкой к максимально возможной $\log_2(Q)$.

Где-то в районе $Q = 14$ выигрыш за счёт увеличения размера алфавита уравновешивается возрастанием вероятности ошибки, и скорость передачи достигает своего максимума.

Как и следовало ожидать, более аккуратный анализ ошибок в метрике Ли даёт несколько лучшие показатели в этом случае.

При дальнейшем увеличении $Q$ скорость начинает падать. И если в метрике Хэмминга уменьшение ярко выражено, в метрике Ли оно не велико, а верхняя граница вообще остаётся почти константной.

Это происходит, так как практически все ошибки при передаче меняют символ на соседний. Соответственно, метрика Ли учитывает оное при анализе, чего Хэмминг сделать не в состоянии.
